% The Enforcement Layer Under Test: body of Section~\ref{sec:system} (the \section heading and its
% label stay in main.tex). Labels defined here: tab:sys-precedents, sec:sys-contracts, eq:sys-in,
% sec:sys-retrieval, sec:sys-complete, sec:sys-compose, sec:sys-identity.
% Sources: the library's retriever.py, filters.py, authz.py, caching.py, events.py, mcp.py,
% server.py, docs/security.md and CHANGELOG.md. Completeness is described as on the
% honest-completeness change (the optional exact_search backend attribute, PR #166), which the
% paper reports; the earlier behavior is stated because Section e2e ran before the change.

\sys{}, the layer we test, is a Python library that implements steps (i) to (vi) of
Section~\ref{sec:prob-setup} between a RAG application, an authorizer and a vector store.
Table~\ref{tab:sys-precedents} separates what has precedent from what we add. The enforcement
mechanism is not ours. We add a single reference semantics for permission filters, which makes the
mechanism testable across stores, and a completeness flag that does not overstate what an
approximate index returned.

\begin{table}[tbp]
\centering
\scriptsize
\newcommand{\sysrag}{\raggedright}
\setlength{\tabcolsep}{4pt}
\begin{tabular}{@{}p{0.21\linewidth}p{0.44\linewidth}p{0.3\linewidth}@{}}
\toprule
Component & Precedent & Role in this paper \\
\midrule
\sysrag Pre-filter: the access predicate pushed into the query
 & \sysrag Query modification~\cite{stonebraker1974access}; query integration in multi-user
   full-text search~\cite{buttcher2005security}; vendor pre-filter
   integrations~\cite{cerbos2026implementing,authzed2026langchainspicedb}
 & \sysrag Mechanism under test \tabularnewline
\sysrag Widening post-filter
 & \sysrag Late binding with growing candidate blocks~\cite{bailey2006secure,mangtani2006documentlevel};
   over-fetching two to three times~\cite{openfga2026rag}
 & \sysrag Mechanism under test; its rounds are a side channel (Section~\ref{sec:side}) \tabularnewline
\sysrag Re-check of every candidate
 & \sysrag Late-bound checks per result~\cite{pang2019zanzibar}; real-time verification of
   pre-filtered documents~\cite{aws2026access}
 & \sysrag Violations as an online count of filter false admits (Proposition~\ref{prop:reduction}) \tabularnewline
\sysrag Failing closed
 & \sysrag Withholding documents on transient errors~\cite{aws2026access}; an error instead of a
   partially filtered result~\cite{microsoft2026querytime}
 & \sysrag Mechanism under test \tabularnewline
\sysrag Truman answers over the readable set
 & \sysrag Truman and non-Truman semantics~\cite{rizvi2004extending}
 & \sysrag The flag $c$ as a partial step toward non-Truman answers \tabularnewline
\sysrag Three-operator filter algebra with one reference semantics
 & \sysrag Per-store filter languages
 & \sysrag \textbf{Contribution}: the oracle for differential testing across 14 stores
   (Section~\ref{sec:conformance}) \tabularnewline
\sysrag Honest completeness
 & \sysrag Exact search when fewer than $k$ approximate results return~\cite{opensearch2026efficient}
 & \sysrag \textbf{Contribution}: an exactness declaration and a flag that never reports a short
   approximate answer as complete (Sections~\ref{sec:sys-complete} and~\ref{sec:repair}) \tabularnewline
\sysrag Composition and tenant scopes
 & \sysrag Policy-combining algorithms~\cite{oasis2013extensible}; denials that override
   grants~\cite{aws2026access}
 & \sysrag Engineering; described, not claimed \tabularnewline
\sysrag Decision cache with a revocation bound
 & \sysrag Snapshot tokens against applying stale ACLs~\cite{pang2019zanzibar}; revocation defeated
   by a cached index~\cite{roychowdhury2024confusedpilot}
 & \sysrag Engineering; described, not claimed \tabularnewline
\sysrag Subject from a verified token
 & \sysrag Token exchange~\cite{jones2020oauth}; MCP authorization~\cite{mcp2026model}
 & \sysrag Adopted, not claimed \tabularnewline
\bottomrule
\end{tabular}
\caption{Components of the enforcement layer, their precedents, and their role in this paper.}
\label{tab:sys-precedents}
\end{table}

\subsection{Contracts and filter semantics}\label{sec:sys-contracts}

For a subject $s$ and an action, an \emph{authorizer} offers $\mathrm{check}(s,H)$, which every
authorizer implements and which returns one Boolean per hit, in order; $\mathrm{filter\_for}(s)$,
which returns a filter $\varphi_s$ selecting \emph{exactly} what $s$ may read, or nothing; and,
optionally, $\mathrm{scope\_for}(s)$, which returns a \emph{necessary} condition $\psi_s$, such as
the subject's tenant, that every readable chunk satisfies. Besides ACLs copied into chunk metadata,
user functions and tenants, \sys{} supports seven engines: OPA (through a filter document its policy
returns), OpenFGA, SpiceDB and Cerbos offer
filters~\cite{openfga2026perform,authzed2026protecting,cerbos2026cerbos}, and Ory Keto, Cedar and
Casbin offer checks only. A \emph{search backend} declares whether
$\mathrm{search}(q,k,\varphi)$ applies a filter inside the search, and returns at most $k$ hits,
best first, with metadata exactly as stored, since the re-check decides on it (A2). Adapters cover
14 stores (pgvector, Qdrant, Milvus, MongoDB Atlas, Weaviate, Redis, LanceDB, FAISS, Vespa,
Cassandra, Pinecone, Elasticsearch, OpenSearch and Chroma), and an in-memory backend that evaluates
the reference semantics is their oracle. One algorithm, written as a generator of operations, runs
under a synchronous and an asynchronous driver, so both share the security logic.

Filters are built from $\mathrm{In}(f,V)$, $\mathrm{And}$ and $\mathrm{Or}$. Let $\mathcal{V}$ be the
strings and integers, Booleans excluded, with strings compared by code point, without case folding
or Unicode normalization, and never equal to integers. One function, \texttt{matches}, defines, for
metadata $m$ and $V \subseteq \mathcal{V}$,
\begin{equation}\label{eq:sys-in}
m \models \mathrm{In}(f,V) \iff
\begin{cases}
\exists x \in m(f):\ x \in \mathcal{V} \wedge x \in V & \text{if } m(f) \text{ is a list},\\
m(f) \in \mathcal{V} \wedge m(f) \in V & \text{otherwise},
\end{cases}
\end{equation}
with And and Or as conjunction and disjunction. Every store translation must agree with it, and it
is the oracle of Section~\ref{sec:conformance}. Each design rule closes a leak. (i)~Empty And and Or
are rejected, since an empty conjunction is vacuously true and would allow everything.
(ii)~Values are strings or integers only: Python equates \texttt{True}, \texttt{1} and \texttt{1.0}
while JSON-based stores compare by type, so floats, nulls, Booleans, nested lists and objects never
match, and a bare string such as \texttt{"12"} is rejected rather than read as \texttt{"1"} and
\texttt{"2"}. (iii)~A missing field never matches, so a chunk without permission metadata is readable
by nobody. (iv)~A filter that can match nothing, such as an In with no values, makes translators
raise \texttt{NothingMatches} and backends return nothing without a query, because Pinecone,
MongoDB, Milvus, LanceDB, Redis and Vespa read a missing or empty filter as ``match everything''.
Parsing rejects unknown operators and nesting deeper than 100 levels.

\subsection{Retrieval}\label{sec:sys-retrieval}

By default \sys{} pre-filters when the backend supports filters and the authorizer returns
$\varphi_s$, and post-filters otherwise, inside $\psi_s$ when one is offered; a caller who demands
pre-filtering gets an error rather than a silent fallback.

\paragraph{Pre-filter and re-check.}
One search with $\varphi_s$ returns the store's answer $H_0$, truncated to $k$ so that a store that
ignores $k$ cannot widen the checks. Every hit is checked again; hits that fail are dropped, not
replaced, and counted as \emph{violations}. Because filter and check are separate code paths, a
violation is, at a single revision, a false admit by the filter or its translation, observed online
(Proposition~\ref{prop:reduction}). With the re-check off, hits are still matched against
$\varphi_s$ locally, so a store that ignores the filter still cannot leak.

\paragraph{Widening post-filter.}
Without a filter, \sys{} fetches $k \cdot o$ candidates ($o = 4$ by default), checks those not seen
before, and multiplies the fetch by $o$ until $k$ hits are permitted, the store returns fewer
candidates than requested, or a cap is reached (1,000 by default; a larger $k$ is refused). A
redactor and a reranker, if configured, receive only permitted hits and cannot add one.

\paragraph{Failing closed.}
Any exception from the authorizer ends the request with an authorization error and no hits. The
message is fixed (``authorizer failed''), the original exception kept only as its cause, because
agent frameworks show tool errors to the model and exception text can name hosts or credentials.
Malformed answers fail likewise: a filter that is not an In, And or Or, a decision list of the wrong
length, or any decision that is not a Boolean (\texttt{1}, \texttt{"no"} and \texttt{None} are
rejected, not coerced).

\subsection{Honest completeness}\label{sec:sys-complete}

Every result carries the flag $c$ of Section~\ref{sec:prob-setup}. Earlier releases set $c = 1$ for
any pre-filtered result without violations, an empty one included, and for any post-filter whose
last search came back short. In the run of Section~\ref{sec:e2e}, \sys{} thus reported empty answers
as complete for least-privileged subjects who could read hundreds of chunks; the project's notes
attribute them to pgvector's iterative scan stopping at its tuple or memory limit, a diagnosis that
Section~\ref{sec:quality} tests. The version we evaluate lets a backend declare
\texttt{exact\_search} when it ranks every chunk that satisfies the filter, so that a short answer
proves nothing else matches; only the in-memory reference and FAISS over a flat index (its default)
do. With $e$ the declaration and $v$ the number of violations, a pre-filtered result has
$c = \bigl(|H_0| = k \vee e\bigr) \wedge (v = 0)$; a post-filtered one has $c = 1$ only if $k$ hits
were permitted, or the last search came back short before the cap and either no scope applied or $e$
holds. The flag is thus honest about under-retrieval, not mis-retrieval: a full answer is taken at
the store's word, so $c = 1$ does not certify maximality (Definition~\ref{def:criteria}), and we
measure mis-retrieval offline (Section~\ref{sec:quality}). An unfiltered post-filter also still reads
a short list as the end of the store, which fails for pgvector with iterative scans off (at most
\texttt{hnsw.ef\_search} rows, 40 by default). Count-based certification and a bounded exact
fallback, as in OpenSearch~\cite{opensearch2026efficient}, are evaluated in Section~\ref{sec:repair}.

\subsection{Composition, tenants and caching}\label{sec:sys-compose}

\paragraph{Composition.}
$\mathrm{AllOf}$ admits a chunk when every part does, asking each part only about chunks that the
earlier parts admitted; $\mathrm{AnyOf}$ asks each part only about chunks still denied. Their
filters, the conjunction or disjunction of the parts' filters, exist only when every part has one.
Their scopes are asymmetric. A scope is valid when $R(s) \subseteq \{d : m_d \models \psi_s\}$.
Since $\bigcap_i R_i \subseteq R_j$ for every $j$, AllOf may conjoin whichever filters or scopes its
parts offer; since leaving a part out of a scope for $\bigcup_i R_i$ can hide readable chunks, AnyOf
offers a scope only when every part offers a filter or a scope. Every candidate is still checked, so
a wrong scope costs speed or recall, never security. The tenant authorizer admits chunks whose tenant
field equals the subject's tenant and raises when the subject has no usable tenant, so a request
that lost its tenant returns nothing rather than everything; combined with an engine that offers no
filter, it becomes the post-filter's scope.

\paragraph{Caching.}
A caching wrapper stores, per subject and action, the filter, the scope, and each decision keyed by
the chunk's identifier and metadata. Its time to live (30\,s by default) is a stated revocation
bound, and an explicit invalidation ends it at once, also for answers that calls in flight obtained
before it. Keys encode type and length, so \texttt{1}, \texttt{"1"}, \texttt{True} and \texttt{1.0}
never share an entry, and failures are never cached. Entries shared through Redis carry their
absolute expiry, so the bound holds across replicas, and, with a signing key, an HMAC-SHA256 over key
and content, so whoever can write to Redis cannot plant a decision. The bound covers the cache alone:
engine staleness~\cite{authzed2026consistency} and the freshness of ACLs copied into chunks add to
it. The cache weakens (A1) of Proposition~\ref{prop:reduction} to a decision at most one time to live
old, and its invalidation needed three successive fixes for races that stored a pre-invalidation
answer after the invalidation.

\subsection{Identity and audit}\label{sec:sys-identity}

No interface lets a client or a model choose the subject. The MCP search tool takes only a query and
a bounded result count, derives the subject from the access token the server verified for the
request, and refuses calls without one~\cite{mcp2026model}. When an agent acts for a user, the token
is issued for the user, for example by token exchange, with the agent as the client or
actor~\cite{jones2020oauth}; the agent then holds no permission its user lacks~\cite{mitre2026atlas},
and no prompt can redirect the tool to another user's permissions, the confused-deputy
pattern~\cite{hardy1988confused}. The HTTP service for other languages verifies OIDC bearer tokens,
each with its key's own asymmetric algorithm, and refuses bodies with unknown fields, so a request
cannot name a user, group or tenant. Every call, failed and cancelled ones included, emits one event
with the subject's identifier, the returned identifiers, the counts, $c$, timings and the error's
type, but never query text, query vectors or chunk text; given a key of at least 16 bytes, it adds an
HMAC-SHA256 of the query, so queries can be correlated without being stored. The violation count and
$c$ are the online signals that Sections~\ref{sec:quality}, \ref{sec:repair} and~\ref{sec:e2e}
compare with offline oracles.
